\documentclass[instructions.tex]{subfiles}
\begin{document}
 
\chapter{Du respect que l’homme doit avoir pour son âme.}
 
Si notre corps, quoique formé de terre et de boue, mérite tant de respect, étant formé à l’image du corps adorable de l’Homme-Dieu, combien plus d’honneur mérite notre âme, qui est formée à l’image de la Divinité! Oui, notre âme est l’image de Dieu : Il fit les animaux et les reptiles de la terre, chacun selon son espèce, dit l’Écriture ; il dit ensuite : Faisons l’homme à notre image et ressemblance\footnote{Et ait : Faciamus hominem ad imaginem et similitudinem nostram. Gen. 1. 26.}.
 
C’est en ce point que les hommes charnels et les incrédules font connaître la bassesse de leurs sentimens. Au lieu de s’élever jusqu’à Dieu pour connaître les rapports et la ressemblance de l’homme avec son Créateur, ils s’appliquent à étudier les rapports qu’ils ont avec les brutes, afin de se persuader qu’ils ont une âme de même espèce que les ours et les pourceaux. Nous laissons ces monstres, ces hommes charnels, qui sont l’opprobre de l’humanité, dans le rang où ils se placent.
 
Rien de plus digne de nos attentions et de nos admirations, que les rapports de notre âme et sa ressemblance avec Dieu, son créateur. Dieu est un esprit pur, intelligent et libre, immortel et infini; notre âme est aussi un esprit : elle a l’intelligence et la liberté; elle est immortelle, elle est comme infinie.
 
\primo Qu’est-ce qu’un esprit? C’est un être pur et simple, sans composition, sans parties, sans matière : c’est quelque chose de plus noble, de plus parfait que tous les astres; tellement au-dessus de nos sens, qu’il nous est aussi impossible d’imaginer la beauté et la perfection d’un esprit, qu’à un aveugle qui n’a jamais vu le jour, d’imaginer et de discerner les couleurs.
 
Un esprit pense, agit, opère et se détermine par lui-même; la matière et le corps ne le peuvent. Les opérations des êtres corporels, des astres, des planètes, des animaux, sont grossières et matérielles, les opérations de l’âme sont au-dessus de la matière; elle peut se connaître, réfléchir sur elle-même; elle rappelle le passé; elle suppute les jours, les années et les siècles; elle va creuser dans l’antiquité, pour s’instruire par les histoires et les événemens; elle pénètre dans l’avenir par des raisonnemens et des conjectures solides.
 
Sans sortir d’elle-même, l’âme parcourt l’univers par sa pensée; elle va, dans un moment, d’un pôle à l’autre, de l’orient à l’occident; elle mesure l’étendue des cieux, le mouvement et la grandeur des astres; elle connaît leurs combinaisons; elle prédit leurs rencontres et leurs éclipses; elle examine, elle pénètre les ressorts de la nature; elle découvre ses secrets, les propriétés des plantes et des minéraux; elle va fouiller jusque dans le centre de la terre, pour en tirer ce qu’il y a de plus riche et de plus utile.
 
Bien plus, elle peut s’élever jusqu’à Dieu, et connaître ses grandeurs; elle en raisonne, elle les adore, et, quoique Dieu habite une lumière inaccessible, elle le découvre par son intelligence, et s’unit à lui par son amour; toutes ces nobles opérations sont au-dessus de tout ce qui n’est pas esprit, et la matière la plus subtile ne pourra jamais y atteindre. Telles sont les opérations de l’âme; elle est donc spirituelle et intelligente.
 
\secundo La liberté est un autre apanage de notre âme; liberté qui la distingue si fort des animaux. C’est en vain qu’on cherche à connaître quelle est la liberté des animaux, quelles sont leurs connaissances, quels principes les font agir; c’est un secret que Dieu a voulu nous cacher, et que nous ne pénétrerons jamais: connoissons ce que nous sommes, et laissons là les bêtes.
 
Un homme qui se persuade que les animaux agissent comme nous, est un fou, et ne comprend pas sa noblesse\footnote{Homo, cùm in honore esset, non intellexit; comparatus est jumentis. Ps. 48. 13.}.
 
Non, les animaux n’ont point la vraie liberté, parce qu’ils ne sont capables ni de raisonner, ni de délibérer; ils agissent toujours selon que l’impression des objets extérieurs les détermine, ou selon l’instinct matériel que Dieu leur a donné. Ils ont partout, chacun dans son espèce, les mêmes opérations, les mêmes cris, les mêmes hurlemens, le même chant, les mêmes façons, et les auront toujours, parce qu’ils ne sont pas libres d’agir autrement.
 
Les hommes, au contraire, dans toutes les contrées de l’univers, ont différentes coutumes, différens sentimens, différens langages, différens accens, différentes manières les uns des autres; preuve qu’ils ont la liberté en partage, qu’ils ont le pouvoir de prendre tel langage, telles façons, telles coutumes qu’il leur plaît.
 
Pourquoi l’homme, malgré son penchant au mal, peut-il se vaincre lui-même, et faire le bien quand il veut? et pourquoi, malgré sa raison qui approuve le bien, malgré la loi de Dieu qui le lui inspire, peut-il faire aussi le mal? C’est parce qu’il est libre, et l’arbitre de sa volonté.
 
C’est pour cela que Dieu, qui voit les crimes, ne les empêche pas, quoiqu’il les défende par sa loi et les punisse par sa justice; parce qu’ayant créé l’homme libre, il ne veut point donner atteinte à notre liberté. Il l’a laissé entre les mains de son propre conseil; il lui présente le feu et l’eau, c’est à lui de choisir.
 
Liberté glorieuse à l’homme! sans cela il n’aurait ni mérite ni récompense; mais liberté funeste quand il en abuse! Plus vous êtes libre et maître de faire tout ce qu’il vous plaît, plus aussi êtes-vous criminel et punissable, si vous vivez dans le désordre. La liberté dont Dieu vous a honoré, vous élève au-dessus des animaux, mais elle vous rend plus méprisable que les insectes, lorsque vous en abusez, parce que les animaux les plus vils ne font jamais rien contre l’ordre du Créateur.
 
\section{Résolutions}
 
\primo Je considérerai mon âme comme étant l’image de Dieu. --
 
\secundo Et je ferai tous mes efforts pour ne pas la déshonorer par des pensées, des affections, des désirs et des actions indignes d’elle. 

\prayer{O mon Dieu! gardez vous-même, par votre grâce, votre image que vous avez formée en moi, et ne permettez pas que je la perde pour l’éternité.}
 
\end{document}
